
本文针对微网多阶段电力合约优化调度问题，构建统一的储能与母线平衡优化模型，依次完成四个子问题的仿真求解，得到如下结果：
问题一求得全天计划购电量【待填入：59482.70】$\mathrm{kWh}$，对应的购电总成本【待填入：35126.95】元；问题二采用F1预测搭配0.9分位安全轨迹控制净负荷预测误差，334天全周期总费用为【待填入：1491.75】万元；问题三在0/6/12/18四个事件节点基于实际SOC重协商合同，将应急购电量降至【待填入：8.75】万$\mathrm{kWh}$，对比冻结合同方案节约成本【待填入：134.68】万元；问题四的可调合同方案相较固定合同方案可节约【待填入：148.89】万元，同时大幅削减高价应急购电需求。

上述结果建立在一系列模型假设之上：微网不允许向外售电、外购电无并网功率上限，储能充放电效率取0.9，采用10分钟离散时间步长仿真。在完成图14、图20--23对应的数值实验后，可以进一步评估不同预报时刻、储能参数变化对调度方案的影响，量化参数对系统经济稳健性的作用。